\chapter{Introduction to UFOLOGY}\label{ch:intro}
\markboth{Introduction to Ufology}{Introduction to Ufology}

\chapquote{``Ufology is more than just flying saucers, and men who wear tin foil hats. It's the study of claims and how one goes about discerning fact from fiction.''}{Evan Svendsen}

% ---------------------------------------------------------------
\section{Redefining The Undefinable}\label{sec:redefiningundefinable}

Before I present some arbitrarily specific definition for the word: UFOLOGY, I
intend to take a step back for a moment. In effort to make it clear to the
audience that the following presented is a definition of ufology as defined by
myself, and myself alone. There have been many countless researchers, and
scientists that have come before me to bear the title of: UFOLOGIST, and they may
interpret the word in an entirely different manner to my own; and that's OK. They
may disregard everything I put forth, however most of the knowledge I have gained
on the subject matter has in one form or another come from my infatuation with the
research of these scientists.

So begin with the shortest version anybody offers. Ufology is the study of
unidentified flying objects. That sentence takes about two
seconds to say and roughly seventy years to argue about, because every single word in
it is doing more work than it appears to.

Start with \emph{unidentified}. This is a statement about the observer, not about the
object. \hi{When I say a thing is unidentified, I am telling you something about the limits
of our identification, and nothing whatsoever about what the thing is.} This is the
single most abused word in the field. The moment somebody says ``it was a UFO,'' and
means by that ``it was a spacecraft,'' they have smuggled a conclusion into a
description. Section~\ref{sec:fallacylanguage} warned you about exactly this trick. The word describes a
failure of our records, and failures of our records are extremely common and not
especially exciting.

Then \emph{flying}. Half the interesting reports involve objects that were not
obviously flying in any aerodynamic sense --- hovering with no visible lift mechanism,
accelerating without a sonic boom, or entering water. ``Flying'' presumes a mode of
travel we recognise, which rather undercuts the ``unidentified'' part.

And \emph{object}. This is the boldest assumption of the three. It asserts that there
is a material thing out there rather than an optical event, a plasma, a reflection, a
radar artefact, or a perceptual phenomenon happening inside a skull. \hi{A great deal of
what gets reported is real without being an object at all.}

Which is why the terminology keeps mutating. The US government now prefers
\term{UAP} --- Unidentified Anomalous Phenomena, having earlier stood for Unidentified
Aerial Phenomena. The change from ``aerial'' to ``anomalous'' was made specifically to
accommodate transmedium and submerged reports. I have mixed feelings about the
relabelling. On the one hand, ``phenomena'' is more honest than ``object'' and drops
the assumption of a craft. On the other hand, ``anomalous'' quietly reintroduces a
different assumption --- that the thing is strange --- when most reports resolve to
something entirely ordinary.

\begin{funfact}
The term ``flying saucer'' is a journalistic accident. In June 1947 pilot Kenneth
Arnold described nine objects near Mount Rainier as moving ``like a saucer if you skip
it across water''. He was describing the \emph{motion}, not the shape --- he said the
objects themselves were crescent-shaped. A reporter compressed this into ``flying
saucers,'' and within weeks the entire world began seeing discs. It is the cleanest
demonstration in the field's history that people frequently see what they have been
told to expect.
\end{funfact}

So here is the working definition I use, and it is deliberately unglamorous. Ufology is
the systematic study of reports of aerial and transmedium phenomena that competent
investigation has failed to identify, together with the physical, historical, and
psychological questions those reports raise. Note what that definition does not contain.
It does not contain aliens. Extraterrestrial visitation is one hypothesis within
Ufology, and by no means the most defensible one. \hi{A field defined by its preferred
answer is not a field; it is an advocacy group with footnotes.}

Stated less formally, and this is the version I would give somebody who stopped me
in the street: Ufology in its simplest of forms is a study of claims. Claims that
are made about the supernatural. Claims about: alien abduction, crop circles, time
travel, life in the universe, and parapsychology. For it CAN ONLY BE study of
claims, as there is no way to objectively recreate on demand for the purposes of
testing in a controlled scientific environment these UFOs. Welcome to the science
of the fringe.

One last piece of vocabulary before we go on, because you will meet it constantly. An
\term{IFO} is an identified flying object --- a report that was investigated and
resolved. IFOs are not failures. IFOs are the denominator. Any honest study of this
subject spends the overwhelming majority of its time producing IFOs, and the small
residue that resists explanation is the only thing that has any claim on a scientist's
attention.

% ---------------------------------------------------------------
\section{The Advent of AI}\label{sec:ai}

I promised in Section~\ref{sec:correlationvscausation} that I would come back to artificial intelligence, and here
we are. The connection is not that computers will solve Ufology, though they may help.
The connection is philosophical: AI is our best available laboratory for studying how
intelligence can appear where nobody put it.

A deep neural network is, stripped of mystique, a very large pile of numbers that get
adjusted until the output stops being wrong. Nobody programs the behaviour. Nobody
writes a rule that says ``cats have pointed ears''. You show the system several million
examples, the weights settle, and something emerges that behaves as though it
understands cats. Ask it why it decided a particular photograph contained a cat and it
cannot tell you, because there is no reason in there --- only a landscape of adjusted
numbers that happens to slope in the right direction.

This matters to us for three reasons.

First, it destroys the intuition that intelligent-looking behaviour requires an
intelligent designer. Recall the bugs and Conway's cells from Section~\ref{sec:correlationvscausation}. When we
encounter a phenomenon that behaves purposefully --- an object that appears to react to
an approaching aircraft, say, or a light that seems to respond to a flashlight signal
--- our instinct is to infer a pilot. AI is a standing reminder that apparent purpose is
cheap. Thermostats appear to want things. So do slime moulds, and immune systems, and
markets.

Second, if we ever do detect an intelligence out there, there is no strong reason to
expect it to be biological. A civilisation capable of crossing interstellar distances
has almost certainly built machines better than itself, and machines do not require
atmospheres, food, sleep, or a return trip. Which means the whole ``occupant'' folklore
of this field --- the greys, the reptilians, the tall Nordics --- may be an artefact of
our inability to imagine a probe. I will return to this in Chapter~\ref{ch:life}.

Third, and most practically: AI has become an evidentiary catastrophe. As I write this,
any competent user can generate a photorealistic video of a metallic disc over their
own neighbourhood in an afternoon, with correct lighting, correct motion blur, and
correct camera shake. The entire photographic corpus of Ufology --- seventy years of
contested stills and grainy footage --- has just been retroactively devalued, because
nothing that arrives from now on can be authenticated by inspection alone.

\begin{funfact}
This cuts the other way too, and almost nobody mentions it. If a genuine image were
captured tomorrow, it would be indistinguishable from the fakes and therefore worthless
as evidence. AI has not merely made faking easier. It has made proof by photograph
structurally impossible. Whatever this field becomes, it will have to become
instrument-based --- radar, spectrometry, calibrated multi-sensor arrays --- or it will
become nothing at all.
\end{funfact}

An adaptive-predictive model does not care what domain you point it at. Trained on the sky it tells
you where to put a camera; trained on genomes it tells you what a mutation will do before anybody
performs it. That second application has quietly dismantled the best argument ever made for a
scientific prohibition --- the argument that some experiments cannot be run because their
consequences cannot be foreseen. Consequences can now be foreseen, imperfectly and improvingly, in
software. I take that up properly in Section~\ref{sec:animalhumanhybridization}, and I mention it
here only so the reader knows that the advent of AI is not merely an evidentiary problem for this
field. It is also the first tool we have ever had that lets a claim about the future be tested
before it is inflicted on anybody.

There is a fourth reason, and it is the one I care about most, because it is the only place
in this book where I have to declare an interest. The four-way classification I teach in the
lecture version of this chapter --- preprogrammed or adaptive, crossed with reactionary or
predictive --- leaves exactly one box that matters to us. Adaptive-predictive AI: a system
that is trained rather than programmed, and whose output is a probability about something
that has not happened yet. The example I used in the lecture hall was \hi{PredPol}, the
predictive-policing company founded in 2012, which I first came across in 2015 in a
documentary presented by the physicist Michio Kaku, and whose promotional video I played to
the class. The point of showing it was not policing. The point was the shape of the claim:
here is a system that says \emph{something will happen there, at roughly that time}, and is
therefore capable of being wrong in public.

\begin{funfact}[Fun Fact: A Prediction on the Record]
\funfactlead{I wrote this down in the autumn of 2017, before any of the hardware existed.}
The original lecture notes for this section end with a paragraph I have never revised, and I
reproduce it here exactly as it was written:

\begin{quote}\itshape\small
``I propose that this idea may be the first ever scientifically structured method for the
investigation of UFOs! It is scientifically based because it solves the problem of
un-falsifiability. We can make a hypothesis with such technology that a UFO will appear over
this city, and it will be that colour and approximate size, and will be seen by 3 people, and
photographed by 2, and 1 person will experience missing time, and then we can actually go AND
TEST if that in fact happens!''
\end{quote}

\noindent That is, near enough, a description of what the UFO Data Acquisition Project became
--- automated multi-sensor stations that trigger on a machine-vision detection, run through
the quiet hours, and produce a record you can be wrong about. See
Section~\ref{sec:ufodapautomatedmultise}. I claim no credit for the engineering; Ron Olch
built it and I did not. But I am going to permit myself the small satisfaction of having
argued for a falsifiable, prediction-first ufology in a student lecture hall a couple of
years before the field went and built one.
\end{funfact}

\begin{funfact}[Fun Fact: This Book Practises What the Chapter Preaches]
Since the previous box just told you that AI has made proof by inspection impossible, honesty
requires the disclosure that follows. \hi{This entire book was written with AI, from scratch.}
Every chapter began as my own lecture notes, my own arguments, my own sceptical positions and
my own bad jokes --- and every chapter was then drafted, structured, expanded, fact-checked
against sources, typeset and indexed with a large language model working from those notes at
my direction. The judgements are mine and I own the errors. The prose is a collaboration.
I mention it for the same reason I mention everything else in this book: a claim you cannot
check the provenance of is worth less than one you can, and it would be absurd to spend four
hundred pages demanding that a field declare its methods and then quietly not declare mine.
\end{funfact}

% ---------------------------------------------------------------
\section{The Most Important Image Ever Taken By Mankind}\label{sec:mostimportantimageeverta}

\begin{center}
\adjustimage{max size={4.8in}{5.4in}}{hubble_deep_field.jpg}\par\vspace{4pt}
{\RaggedRight\scriptsize\color{steel}\textbf{\textcolor{ufogreenink}{The Hubble Deep Field.}}
Hubble Space Telescope, WFPC2, 18--28 December 1995. 342 exposures, about 140 hours,
covering 2.6 arcminutes of Ursa Major --- roughly one twenty-four-millionth of the sky.
Nearly every object in this frame is a galaxy. Photograph: R.~Williams and the Hubble
Deep Field Team (STScI) / NASA, public domain.\par
\vspace{1pt}{\color{midgrey} Image sourced from:} \srclink{https://esahubble.org/images/opo9601a/}\par}
\end{center}
\vspace{6pt}

Hold your arm out at full stretch and hold up a grain of sand. The patch of sky that
grain covers is roughly what the Hubble Space Telescope stared at for ten days in
December 1995. There was nothing there. That was the entire point. And what came back is,
I will argue for the rest of this section, the most important photograph our species has
ever taken.

Start with the arithmetic, because the arithmetic is the argument. The field of view was
2.6 arcminutes on a side. An arcminute is one sixtieth of a degree; the full Moon is
about thirty arcminutes across. So the target was a square smaller than a tenth of the
Moon's face, which works out to roughly one twenty-four-millionth of the whole sky ---
about the angular size of a tennis ball a hundred metres away. Hubble pointed at that
speck, in an unremarkable stretch of Ursa Major just off the handle of the Big Dipper,
and held still.

It was chosen precisely because it was boring. Robert Williams, then Director of the
Space Telescope Science Institute, wanted a patch of sky with nothing in it: no bright
foreground stars to bloom across the detector, no nearby galaxies, no dust from our own
galactic plane in the way, and a location Hubble could keep in view continuously without
the Earth swinging into the frame. Then he spent his Director's Discretionary
observation time --- the slice of the schedule the Director may allocate personally, the
most jealously watched commodity in astronomy --- on ten straight days of photographing
nothing. Colleagues told him it was a career-ending waste of the most oversubscribed
instrument in the world.

You have to understand what the telescope's reputation was at that point. Hubble
launched on 24 April 1990 to enormous fanfare, and on 27 June 1990 NASA announced that
its primary mirror was the wrong shape. Not badly wrong --- exquisitely, catastrophically
slightly wrong. The 2.4-metre mirror had been ground too flat towards its edge, by about
two micrometres, roughly a fiftieth of the width of a human hair. That is
spherical aberration: light from the middle of the mirror and light from the rim refuse
to meet at the same focal point, so every star arrives as a small fuzzy halo instead of a
point. The cause, traced later, was mundane and awful. A test instrument called a null
corrector, used to verify the mirror's figure during polishing, had one of its lenses
sitting 1.3 millimetres out of position. The mirror was then polished, superbly, to the
wrong prescription. The final error was about ten times the tolerance the design allowed.

A one-and-a-half-billion-dollar telescope became an international punchline. Then, in
December 1993, the crew of STS-61 flew up on \emph{Endeavour} and, over five spacewalks,
installed a new camera with the inverse error ground into its own optics along with a
unit called COSTAR, which put small corrective mirrors in front of the other
instruments. In effect they fitted the telescope with spectacles. The optics worked, the
images sharpened, and the obvious next question was: how good is it \emph{really}?

The Deep Field was the answer to that question. It was not primarily a discovery mission.
It was a test --- the hardest test anyone could design for a repaired telescope. Take the
faintest, emptiest thing you can find, integrate for as long as you dare, and see whether
the machine can dig a signal out of nothing. Between 18 and 28 December 1995, Hubble
exposed 342 separate frames across about 150 orbits, through four filters spanning the
near-ultraviolet to the near-infrared, roughly 140 hours of shutter-open time in total.
The individual pointings were nudged slightly between exposures so that the software
could later reconstruct detail finer than a single pixel, a technique called drizzling.

What came back was not empty. It was crowded. About three thousand galaxies appeared in
that grain of sand --- spirals, ellipticals, mergers, deformed blue smudges of things
that had not yet decided what shape to be, some of them radiating light that had been
travelling for well over ten billion years before it hit that detector. Fewer than twenty
of the objects in the frame are foreground stars from our own galaxy. Essentially every
point of light you can see is an entire galaxy of hundreds of billions of suns.

Now multiply. If one twenty-four-millionth of the sky holds three thousand galaxies, the
whole sky holds on the order of a hundred billion of them, and later, deeper exposures
pushed that estimate higher still. We aimed a camera at the emptiest thing we could find
and found everything. There is no other photograph in human history that carries that
sentence.

I call it the most important image ever taken because of what it does to the person
looking at it, and it is not what people assume. It is usually filed under humility, and
it does deliver humility, but humility is not the useful part. The useful part is
\emph{scale calibration}, and this field desperately needs it.

The calibration cuts in both directions, which is exactly why I like it. For the
believer, it is cold water. Every one of those smudges contains hundreds of billions of
stars, and the nearest star to our own is still four light years away --- a gap that our
fastest object has not begun to close. Abundance of worlds is not the same as ease of
travel between them; the Deep Field proves the first and says nothing whatsoever about
the second. For the sceptic, it is equally cold water. The volume of space we have
meaningfully examined for other intelligence, set against that photograph, is roughly
equivalent to sampling one glass of water from the Pacific and announcing there are no
fish.

So keep the picture in mind for the remaining chapters. Whenever you meet an argument
that begins ``surely someone would have noticed'' or ``surely they would have come by
now'' --- mine included --- check it against a grain of sand holding three thousand
galaxies, and ask how much of the picture the arguer has actually looked at.

% ---------------------------------------------------------------
\section{SETI}\label{sec:seti}

The Search for Extraterrestrial Intelligence is the respectable cousin of Ufology, and
the relationship between the two is best described as mutually embarrassed. SETI
researchers publish in refereed journals, hold university appointments, and would
rather you did not mention flying saucers in their presence. Ufologists, in turn, regard
SETI as a well-funded exercise in listening to the wrong thing in the wrong place with
the wrong assumptions. Both positions contain some truth.

\figwrap[l]{0.38}{The 26-metre dish at Green Bank used for Project Ozma, the first modern SETI search.}{https://commons.wikimedia.org/wiki/File:Tatel_Telescope.jpg}{fig:ozma}
The basic method has barely changed since Frank Drake pointed a radio telescope at Tau
Ceti and Epsilon Eridani in 1960 and called it Project Ozma (see Figure~\ref{fig:ozma}). You select target stars,
you listen in the radio spectrum --- traditionally near the 1420\,MHz hydrogen line, on
the reasoning that any physicist anywhere would recognise it as a natural landmark ---
and you look for a narrowband signal that nature does not produce. Narrowband is the key
word. Natural sources are broad and messy. A signal squeezed into a few hertz is the
signature of a transmitter.

Sixty-five years of this have produced exactly one candidate anybody still talks about.
On 15 August 1977 the Big Ear telescope at Ohio State recorded a 72-second narrowband
signal near the hydrogen line, close to the expected profile in every respect. Astronomer
Jerry Ehman circled the printout and wrote ``Wow!'' beside it. It has never recurred,
despite extensive re-observation of that patch of sky. One event, unrepeated, is not a
detection. It is a story.

\begin{funfact}
SETI's real problem is not sensitivity, it is arithmetic. Earth's own radio
leakage --- the unintentional stuff, television and radar --- is detectable at
interstellar distances only with instruments far better than any we possess, and our
leakage has been \emph{decreasing} for decades as we move to fibre, satellites and
low-power digital transmission. If our own civilisation is any guide, a technological
species is a bright radio source for perhaps a century and then goes quiet. We may be
listening for a phase almost nobody spends any time in.
\end{funfact}

The deeper criticism, and the one I think lands, is that SETI encodes a specific and
questionable hypothesis: that contact will arrive as a deliberate radio broadcast from a
civilisation that wants to be found. That is a narrow bet. It excludes probes, it
excludes optical and infrared signalling, it excludes technosignatures such as
atmospheric industrial pollutants or waste heat, and it excludes the possibility that
anybody nearby has already been here --- which is precisely the possibility this book
exists to examine.

To be fair, the field has broadened considerably. Optical SETI looks for nanosecond laser
pulses. Breakthrough Listen, funded privately from 2015, has surveyed far more targets
than all previous efforts combined. There is serious work on spectroscopic biosignatures
in exoplanet atmospheres, which sidesteps intelligence altogether and simply asks whether
anything is alive. That is a better question, and we will get to it.

I raise SETI here, this early, for a methodological reason rather than an astronomical
one. SETI is what Ufology would look like if it had grown up properly. It states a
hypothesis, specifies in advance what would count as a detection, publishes its null
results, and reports its candidate events with the caveats attached. It has found nothing
in sixty-five years and it says so plainly. Compare that to a field where a single
unrepeated event becomes a permanent article of faith. The problems with SETI are real,
and I will take them apart properly in Section~\ref{sec:setiproblem}. Its methods,
though, are not the problem. Its methods are the standard we have failed to meet.

% ---------------------------------------------------------------
\section{The Ufologist}\label{sec:ufologist}

What then does it mean to be a Ufologist? What does it take? If no university or
scholarly institute will even indulge in subjects such as this? Where might a person
learn to become a scholarly prophet of such knowledge? To answer these questions, it
must first be stated that ufology involves the professions of nearly every conceivable
job on earth. In what regard you may ask? A Ufologist is anyone who makes considerable
effort of their time to either research, contribute to, or profess the subject. No not
in the fashion of a cult, but in the sheer endless wonder of human curiosity.

Passion is really all it takes. Think of other things in life for instance that one may
be passionate about. If you play an instrument for example, I am sure you have a passion
for your music. But passion alone does not make you a musician, you must also have the
skill of playing said instrument. What about a doctor? Because you passionately watch
videos of open-heart surgery on YouTube? I think I'll take my chances with the cancer.
So, you're not a doctor, and I'm afraid you aren't The Doctor either! What about
Ufologist though? In terms of the skills; it's whatever you bring to the table, and that
skill is driven by passion alone when applied to the field of UFO investigation.

As a result, nearly everyone from physicists to journalists, politicians, historians,
statisticians, philosophers, artists, linguists, authors, archeologists, chemists,
biologists, astronomers, psychologists, theologians, police, pilots, and the military
have in some way either been directly or independently involved in the vast history of
investigating and reporting or professing on the subject of UFOs.

What, then, does a Ufologist actually do? Not what television suggests. There is very
little night-vision footage and a great deal of paperwork.

The honest answer is that a Ufologist is an archivist who occasionally gets to do
fieldwork. The bulk of the job is reading: case files, weather records, air traffic
logs, satellite pass predictions, astronomical almanacs, and the previous forty years of
contradictory literature about whichever case you happen to be looking at. When a report
comes in, the work is not to determine whether it was a spacecraft. The work is to
establish, in order: what the witness actually experienced, what the sky actually
contained at that time and place, and whether the first can be accounted for by the
second.

That sequence matters and it is usually done backwards. Most investigators start with a
hypothesis and go looking. The correct order is to reconstruct the observing conditions
before you form any opinion at all --- because once you have an opinion, Section~\ref{sec:confirmationbias}
tells you exactly what will happen to your subsequent reading of the evidence.

There are, broadly, three species of person in this field, and you should learn to tell
them apart at a distance.

The \textbf{advocate} has already reached a conclusion and is now assembling support for
it. He is often knowledgeable, frequently charming, and completely useless as a source,
because you cannot tell which of his facts survived scrutiny and which merely survived
selection. The \textbf{debunker} is the same animal facing the other way, and is treated
with more respect because his conclusion is the socially safe one. Then there is the
\textbf{investigator}, who is boring. She writes up cases that resolved. She says ``I
don't know'' in print. She publishes the observing conditions even when they undermine
her own earlier work. There are perhaps a few dozen of these people alive and most of
you have never heard of any of them.

What follows is a working guide to the people whose names you will meet again and again,
living first and then the dead, with my own assessment of each attached. I have tried to
be fair and I have not tried to be flattering. Sort them into the three species yourself
as you read; you will find that most of them refuse to sit still in one box, which is
rather the point.

\ufobioportraits

\subsection{Brandon Fugal}

Fugal owns the ground on which the field's most heavily instrumented investigation is
being conducted, and that single fact makes him one of the most consequential figures in
contemporary Ufology despite his having no scientific training whatsoever. Born in
Pleasant Grove, Utah, on 3~April 1973, he built a career in commercial property and is
chairman of Colliers International's Intermountain operation, which is to say he is a
successful broker of office space rather than an investigator of anomalous phenomena. In
2016 he purchased the 512-acre Skinwalker Ranch from Robert Bigelow through a holding
company, remained anonymous as its owner for some three years, and then went public and
became executive producer and on-camera presence of \emph{The Secret of Skinwalker
Ranch}. He is, by his own account, a sceptic who bought a ranch to disprove its
reputation and was changed by what happened there.

My assessment is mixed and I want to be precise about which half is which. To Fugal's
credit, he has spent his own money on instrumentation, has hired people with real
technical credentials rather than psychics, and has kept the property under controlled
access so that the site is not trampled --- all of which is more than most owners of
famously haunted land ever bother to do. Against that stands the structural problem
which no amount of goodwill dissolves: he is simultaneously the funder of the research,
the owner of the data, and a commercial beneficiary of the conclusion that the ranch is
anomalous. Those three roles cannot be occupied honestly at once, not because Fugal is
dishonest but because no one could be trusted in that position, including you and
including me. Add the non-disclosure agreements, the unreleased sensor archives, and the
fact that the findings reach the public through a television edit rather than a
manuscript, and the sceptical verdict writes itself. Until the raw data goes to an
outside laboratory with no financial interest in the answer, the Skinwalker material has
to be classed as unverified. I treat the site itself, and what has and has not been
demonstrated there, at length in Section~\ref{sec:skinwalker}.

\subsection{Travis S.~Taylor}

\portrait[r]{0.36}{uf_taylor.jpg}{\textbf{\textcolor{ufogreenink}{Travis S.~Taylor.}} Photograph: US~Army, public domain.}{https://commons.wikimedia.org/wiki/File:Travis_S_Taylor.png}
Taylor is the closest thing the field has to a genuinely credentialled
physical scientist, and he did not arrive by way of Ufology at all. Born in 1968, he holds
doctorates in optical science and engineering and in aerospace systems engineering, is a
licensed professional engineer, and spent a working career on defence and missile-defence
programmes before joining Radiance Technologies as its chief scientist in 2022. Television
found him through \emph{Rocket City Rednecks} and then installed him as the on-camera
scientist of the Skinwalker Ranch series, which is where most readers will know him. His
virtue is that he actually deploys instruments and actually reports numbers. His problem is
the format he works in: a television season needs anomalies on a schedule, and a
measurement announced to camera before anybody has ruled out the mundane is not a
measurement, it is a cliffhanger. Take the man's physics seriously and treat the editing
with suspicion.

\subsection{Erik Bard}

Bard is the quieter half of the same operation and, to my mind, the more
interesting one. An X-ray physicist by training, he serves as principal investigator at
Skinwalker Ranch and is the person who designs and maintains the sensor coverage --- the
magnetometers, the radiation monitors, the long-baseline arrays. He is markedly more
cautious on camera than the format wants him to be, which is a good sign in a man. The
critique here is not of Bard but of what surrounds him: the ranch has produced years of
instrumented data and very little of it has been released in a form an outside physicist
could reanalyse. Until it is, the work has to be classed as unverified rather than wrong.

\subsection{Linda Moulton Howe}

\portrait[l]{0.36}{uf_howe.jpg}{\textbf{\textcolor{ufogreenink}{Linda Moulton Howe}} with her 1981 Emmy for \emph{A Strange Harvest}. Photograph: Larry W.~Howe, CC~BY~4.0.}{https://commons.wikimedia.org/wiki/File:Linda1981EmmyStrangeHarvest300dpi.jpg}
Howe is the field's most prolific investigative journalist and its
most polarising. Born in Boise in 1942, she took a master's degree from Stanford in 1968,
ran documentary programming at KMGH in Denver from 1978 to 1983, and won an Emmy for
\emph{A Strange Harvest} (1980), the film that turned cattle mutilation into a national
subject. She has run \emph{Earthfiles} for three decades since. I devote a full treatment
to her career and to what her archive is and is not worth in
Chapter~\ref{ch:earthfiles}, so I will only register the summary judgement here: her
documentation of physical cases is unusually diligent and her handling of anonymous
sources is not. When she has a rancher, a vet and a necropsy she is excellent. When she has
a man who cannot be named describing a document nobody else has seen, she is exactly as
reliable as he is.

\subsection{George Knapp}

\portrait[r]{0.40}{uf_knapp.jpg}{\textbf{\textcolor{ufogreenink}{George Knapp}} at the Peabody Awards. Photograph: Peabody Awards / Anders Krusberg, CC~BY~2.0.}{https://commons.wikimedia.org/wiki/File:George_Knapp_(8221331696)_(cropped).jpg}
Knapp is a working newsman who happened to inherit the biggest story in
the field. As chief investigative reporter at KLAS-TV in Las Vegas he has collected
Peabody and Edward R.~Murrow awards for reporting that has nothing to do with UFOs, and it
is worth knowing that before you dismiss him. He broke Bob Lazar in 1989
(Section~\ref{sec:s4lazar}) and co-wrote \emph{Hunt for the Skinwalker}, the book that
created the ranch as a public object (Section~\ref{sec:skinwalker}). Knapp's strength is
that he keeps his sources and keeps coming back to them for decades; his weakness is that
he has never fully separated himself from a story he helped create. He gave Lazar a
national platform without independently verifying the biographical claims, and thirty-five
years of subsequent coverage has been shaped by that first decision. A reporter who is
also a participant is a reporter with a stake.

\subsection{Seth Shostak}

\portrait[l]{0.40}{uf_shostak.jpg}{\textbf{\textcolor{ufogreenink}{Seth Shostak.}} Photograph: NASA / Bowman, public domain.}{https://commons.wikimedia.org/wiki/File:Seth_Shostak.jpg}
Shostak is the field's designated adversary and its most useful one. Born on
20~July 1943, with a physics degree from Princeton and a doctorate in astronomy from
Caltech, he is senior astronomer at the SETI Institute and has spent thirty years
explaining, patiently and in public, why he thinks UFO reports are not evidence of visitors.
His positive case is strong: he works on the one search programme in this territory that
states its methods and publishes its nulls. His negative case is weaker than it sounds,
because it rests on an argument from expectation --- that a real interstellar civilisation
would behave the way a radio astronomer would find legible. That is an assumption, not a
finding. Read Shostak for method and discount him slightly on motive.

\subsection{William J.~Birnes}

Birnes is the field's impresario. He took a doctorate and a law degree
at New York University in 1974, built a career in publishing, bought and edited
\emph{UFO Magazine}, and fronted \emph{UFO Hunters} for three seasons. He is also the
co-author, with Philip J.~Corso, of \emph{The Day After Roswell} --- a book that made the
reverse-engineering claim mainstream on the strength of a single elderly officer's
recollections and effectively no documents. Birnes is genuinely learned about the field's
literature and genuinely uncritical about its best-selling stories. He belongs, without
much argument, in the advocate column.

\subsection{Leslie Kean}

\figwrap[l]{0.36}{No free photograph of Kean exists, so here is her subject instead:
a frame from the 2015 GIMBAL ATFLIR recording, one of the Navy videos her \emph{New York
Times} story of December 2017 put in front of the public.}{https://commons.wikimedia.org/wiki/File:Gimbal_The_First_Official_UAP_Footage_from_the_USG_for_Public_Release.webm}{fig:gimbal}
Kean did the single most consequential piece of work in modern Ufology and
is rarely given credit for it, because the credit went to a newspaper. A journalist rather
than a believer, she published \emph{UFOs: Generals, Pilots and Government Officials Go on
the Record} in 2010 --- a book built almost entirely on named witnesses and released
documents, which is why it still holds up --- and then co-wrote the \emph{New York Times}
story of December 2017 that revealed AATIP and put the Navy videos in front of the public.
Her later book \emph{Surviving Death} (2017) moved into mediums and reincarnation and is a
different kind of exercise entirely, which some readers hold against her. I hold her
method against nobody: she names people and she prints the paperwork. That she also finds
the afterlife interesting tells you about her curiosity, not about her sourcing.

\subsection{Steven Greer}

\portrait[r]{0.36}{uf_greer.jpg}{\textbf{\textcolor{ufogreenink}{Steven Greer.}} Photograph: ``Energy Artist'', public domain.}{https://commons.wikimedia.org/wiki/File:StevenGreer01_(cropped).jpg}
Greer, born in 1955, is the most successful organiser the field has
produced and the one I trust least. A former emergency physician, he founded CSETI and
staged the Disclosure Project event in Washington in 2001, at which more than twenty
military and government witnesses testified on the record --- a real achievement that
deserves to be remembered. He also invented CE-5, the practice of summoning craft by
meditation, sells the training, and has repeatedly announced imminent revelations that did
not arrive. The pattern is the difficulty: the verifiable half of his work is used as the
credential for the unverifiable half. Greer is the clearest illustration in this book of
why an advocate can be simultaneously important and unusable.

\subsection{Roger Leir}\label{sec:leir}

\portrait[l]{0.36}{uf_leir.jpg}{\textbf{\textcolor{ufogreenink}{Roger K.~Leir.}} Photograph: ChurchOfThought, CC~BY-SA~4.0.}{https://commons.wikimedia.org/wiki/File:Roger_Krevin_Leir.jpg}
With the dead the judgement comes easier, because their records are closed.
Roger Leir (20~March 1935 --- 14~March 2014) was a podiatric surgeon who performed
some seventeen operations to remove what he described as alien implants from people who
believed they had been abducted. He documented the surgeries, commissioned laboratory
analyses, and reported anomalous metallurgy and an absence of inflammatory response around
the objects. This is, on its face, exactly the kind of physical-evidence programme the
field needs, and it failed for exactly the reasons such programmes usually fail: no
independent surgical oversight, no blinded pathology, no chain of custody worth the name,
and analyses commissioned rather than replicated. Nothing he removed has ever been shown
by a disinterested laboratory to be anything other than terrestrial. That is the standard
objection and it is the wrong one. Everything he removed from his patients came from inside
the body of a human being. Seventeen implants in total. Seventeen patients. That proves one
thing and one thing alone: the body did not reject them. Why? That is the question this
case leaves my reader with, and it is the question the field has never answered.

\subsection{Stanton Friedman}

\portrait[r]{0.33}{uf_friedman.jpg}{\textbf{\textcolor{ufogreenink}{Stanton T.~Friedman}}, Big Bear, California, 2019. Photograph: SMG2019, CC~BY-SA~4.0.}{https://commons.wikimedia.org/wiki/File:Stanton_T._Friedman_at_Alien_Snowfest_2019_in_Big_Bear,_CA.jpg}
Friedman (29~July 1934 --- 13~May 2019) was the first professional to make
Roswell his life's work, and he was a professional: a nuclear physicist with bachelor's
and master's degrees from the University of Chicago, where Carl Sagan was a classmate, and
fourteen years on nuclear aircraft and rocket programmes before he went full-time in 1970.
He found Jesse Marcel, he co-authored \emph{Crash at Corona}, and he spent five decades
debating anyone who would stand still. Friedman's discipline about documents was real and
his discipline about conclusions was not: he treated the extraterrestrial explanation as
established and everything after that as detail work. He also defended the MJ-12 papers
long past the point at which their typography had convicted them. He was the best advocate
the field ever had, which is a compliment and a limit.

\subsection{Nick Pope}

\portrait[l]{0.36}{uf_pope.jpg}{\textbf{\textcolor{ufogreenink}{Nick Pope.}} Photograph: David Howard, CC~BY~2.0.}{https://commons.wikimedia.org/wiki/File:Nick_Pope_(journalist).png}
Pope (1965 --- 6~April 2026) ran the British Ministry of Defence's UFO desk
from 1991 to 1994 and then spent thirty years as the field's most quotable civil servant,
dying of cancer at sixty. His first book, \emph{Open Skies, Closed Minds}, remains the best
account of what a government UFO office actually does with its post, which is mostly to
ask the radar people and then write a polite letter. He estimated a residual of about five
per cent --- cases the desk could not resolve --- and he was the principal public voice on
Rendlesham Forest (Section~\ref{sec:rendlesham}). The criticism levelled at him is that
his job title inflated in the retelling, from desk officer to something closer to
investigator-in-chief, and that his years on \emph{Ancient Aliens} lent his civil-service
credibility to material that did not deserve it. Both criticisms are fair. His account of
the paperwork is still the most valuable first-hand testimony we have about how states
handle this subject.

\subsection{Frank Drake}

\portrait[r]{0.36}{uf_drake.jpg}{\textbf{\textcolor{ufogreenink}{Frank Drake}} at Cornell, October 2017. Photograph: Amalex5, CC~BY-SA~4.0.}{https://commons.wikimedia.org/wiki/File:Frank_Drake_at_Cornell,_October_2017_(cropped_2).jpg}
Drake (28~May 1930 --- 2~September 2022) was not a Ufologist and would have
been irritated to be listed among them. He pointed a radio telescope at two stars in 1960,
wrote the equation that bears his name on a blackboard at Green Bank in 1961, and thereby
gave the question of other minds a form that could be argued about numerically. I treat the
equation properly in Section~\ref{sec:drake}, including what is wrong with it. He is here
because he is the counter-example that indicts the rest of the field: faced with the same
question, he built an instrument, defined a detection in advance, and published nothing.
That is what nothing looks like when it is reported honestly.

\subsection{Jacques Vall\'ee}

\portrait[r]{0.30}{fg_vallee.jpg}{\textbf{\textcolor{ufogreenink}{Jacques Vall\'ee}} in 2024. Photograph: Christopher Michel, CC~BY-SA~4.0.}{https://commons.wikimedia.org/wiki/File:Jacques_Vall\%C3\%A9e_by_Christopher_Michel_in_2024_02.jpg}
Vall\'ee (born 24~September 1939) is the most intellectually serious person the subject has
produced, which is not the same as being right. Trained in astrophysics at Lille and Paris and
then in computer science at Northwestern, he worked on Mars mapping for NASA, helped build the
ARPANET, and spent the second half of his career as a venture capitalist --- a curriculum vitae
that makes the usual complaint about credulity difficult to sustain. He came to the subject as
Hynek's collaborator and left it as Hynek's critic.

His contribution was to notice, in \emph{Passport to Magonia} (1969), that the reports do not
look like visits and do read remarkably like the fairy and demon encounters of earlier centuries
--- same abductions, same missing time, same lights, same absurd little details --- which is a
genuine observation about the data and one the nuts-and-bolts school has never answered. What he
built on it, the \term{control system} hypothesis, holds that the phenomenon is a feedback
mechanism conditioning human belief. I think the observation is sound and the hypothesis is
unfalsifiable: a control system that adapts its manifestations to the beliefs it is shaping
cannot be caught out by any measurement, and a proposition no measurement can embarrass is not a
scientific one. He is here because he asked the best question in the field's history and then
answered it with something no instrument can test.

\subsection{J.~Allen Hynek}

\portrait[l]{0.27}{uf_hynek.jpg}{\textbf{\textcolor{ufogreenink}{J.~Allen Hynek.}} Photograph: US~Government, public domain.}{https://commons.wikimedia.org/wiki/File:Allen_Hynek_Jacques_Vallee_1.jpg}
Hynek (1~May 1910 --- 27~April 1986) is where any list like this has to
end. An astronomer at Ohio State and later Northwestern, he took the Air Force's shilling
in 1948 to explain sightings away, did it well, coined the phrase that dogged him for life,
and then --- slowly, over twenty years of Blue Book case files --- changed his mind. He
founded the Center for UFO Studies in 1973, gave us the close-encounter scale, and
consulted on the Spielberg film that took its title from it. He died of a brain tumour in
Scottsdale at seventy-five. His failing was that his conversion eventually outran his
rigour: the later Hynek entertained interdimensional hypotheses on evidence he would have
rejected in 1955. That does not diminish him. He is the only major figure in this field who
demonstrably followed the evidence somewhere he did not want to go.

\ufowrapclear
\begin{funfact}
The most consequential Ufologist in history did not believe in UFOs when he took the job.
Dr.~J.~Allen Hynek was hired by the US Air Force in 1948 as an astronomical consultant
specifically to explain sightings away, and he did so cheerfully for years --- he is the
man who coined the ``swamp gas'' explanation that got him mocked for the rest of his
life. It was the sheer volume of residual cases crossing his desk that eventually turned
him, not any single sighting. He went on to devise the close-encounter classification
system we still use. A man who changed his mind after twenty years of evidence is worth a
thousand who were certain from the start.
\end{funfact}

\ufowrapclear
\ufostdportraits

\subsection{Credentials, Including Mine}

Which brings me to credentials, and to my own. There is no degree in Ufology, as I said
in the Introduction. This is a genuine structural problem and not merely a personal
inconvenience, because it means the field has no mechanism for excluding incompetence and
no mechanism for accrediting competence. Anybody can print a business card. I did. What
substitutes for accreditation is your paper trail: the cases you have written up, the
sources you have traced, and --- most of all --- the number of times you have publicly
concluded that something exciting was in fact Venus.

\begin{srcnotes}
\item Biographical details for the living figures are drawn from their own published
bios, employer pages and interviews: Radiance Technologies (Taylor), the SETI Institute
(Shostak), KLAS-TV (Knapp), \emph{Earthfiles} (Howe) and CSETI (Greer).
\item Kean, L.\ (2010). \emph{UFOs: Generals, Pilots and Government Officials Go on the
Record}. New York: Harmony; and Cooper, H., Blumenthal, R.\ \& Kean, L., ``Glowing Auras
and `Black Money': The Pentagon's Mysterious U.F.O.\ Program'', \emph{New York Times},
16~December 2017.
\item Pope, N.\ (1996). \emph{Open Skies, Closed Minds}. London: Simon \& Schuster.
\item Friedman, S.\ \& Berliner, D.\ (1992). \emph{Crash at Corona}. New York: Paragon
House.
\item Hynek dates and cause of death per the \emph{New York Times} obituary, 1~May 1986,
and the J.~Allen Hynek Papers finding aid, Northwestern University Archives.
\item Drake dates per University of California in memoriam notice and contemporaneous
obituaries, September 2022.
\end{srcnotes}

% ---------------------------------------------------------------
\section{UFOs Around the World}\label{sec:ufosaroundworld}

One of the more useful things you can do with a global dataset is look for the places
where the reports \emph{differ}, because that tells you which features belong to the
phenomenon and which belong to us.

The American pattern is the one everybody knows: discs, then triangles, then a heavy
emphasis on government concealment and on abduction. The abduction narrative in
particular is startlingly American. It arrives in the early 1960s, it involves medical
examination and missing time, and it maps almost too neatly onto the anxieties of a
society with a new relationship to hospitals and to surveillance. We will take that apart
properly in Chapter~\ref{ch:abduction}.

Brazil has the highest concentration of physical-effect and injury cases anywhere in the
world, and it has the one case I cannot file away tidily.

\subsubsection*{Varginha, Brazil, 20 January 1996}

Varginha is a coffee town of about 135,000 people in Minas Gerais, four hundred kilometres
inland from Rio. On the afternoon of 20~January 1996, in the middle of a violent summer
storm, three young women --- the sisters Liliane and Valqu\'iria Silva and their friend
K\'atia Andrade, then aged between fourteen and twenty-two --- crossed a vacant lot on
their way home and stopped dead. Crouched against a wall in the weeds was a being roughly
four feet tall: brown, oily skin, an oversized head with three bumps on the forehead,
enormous red eyes, and a smell they could not place. They said it looked frightened. They
ran, and told their mother they had seen the devil.

That is where most retellings begin and where the sceptical account normally ends. A 2010
Army inquiry --- two volumes, some six hundred pages, digitised and now public --- concluded
that the three had misread a local man in the rain: Lu\'is Ant\^onio de Paula, known in
Varginha as \emph{Mudinho}, a resident with mental disabilities who habitually walked in a
crouch. Brian Dunning's treatment on \emph{Skeptoid} is representative of the tone of the
wider dismissal: a town with a tourist industry, a story that grew in the telling, no
physical evidence worth the name. All of that is fair. Both surviving witnesses still
reject the Mudinho identification --- they had known the man since childhood --- but witness
confidence is not evidence, and I would not build a paragraph on it.

Varginha is also the case that brought the film-maker James Fox to Brazil. His
\emph{Moment of Contact} (2022) is a full-length treatment of these three days, built the way
all of his work is built --- witnesses on camera, in their own words, with the questions left
in --- and it is the reason most English-speaking readers have heard of the case at all.
Chapter~\ref{ch:media} takes up what Fox's seven films have and have not accomplished, and it
returns to this one, because \emph{Moment of Contact} is where his method shows both of its
edges most clearly: it preserved testimony that would otherwise have died with the witnesses,
and it also poured money and pilgrims into a town that already had a saucer on its water
tower. I use it here for the interviews and nothing else. A documentary is a collection of
witnesses, edited; it is not an investigation, and Fox has never claimed it was.

What I cannot get past is the funeral.

\actualq[question]{Did three teenagers really see an extraterrestrial in a vacant lot in
Minas Gerais?}{Do hoaxes normally end in a \hi{death}?}

Military Police officer Marco Eli Chereze was twenty-two. On the evening of 20~January he
was reportedly part of the detail that handled the creature --- witnesses said he took hold
of it with his bare hands and put it in a vehicle. On 7~February he was admitted to
hospital with severe pain and loss of function in one arm. He was operated on for an
abscess in the armpit. He died on 15~February. The cause of death on the record is
generalised infection --- sepsis, with pneumonia --- from a bacterium that, by the account of
the doctors involved, had no business killing a healthy young man in a fortnight. His
former battalion commander, Maur\'icio Ant\^onio Santos, has stated flatly that the death was
a post-operative hospital infection and that Chereze was involved in no incident with
anything extraterrestrial. That is a coherent explanation, and hospital-acquired sepsis
kills healthy young people every week of the year in every country on Earth.

So the death is not proof of anything. But notice what it does to the shape of the claim.
A hoax is a story: it costs nothing, it can be extended indefinitely, and its authors
typically stay alive to enjoy it. Roswell produced a weather balloon and a small industry.
The Belgian wave produced a confessed fake photograph. Varginha produced a grave. When a
fabrication has a body attached to it, one of three things is true --- the death is pure
coincidence, the death is being retro-fitted to the story, or the story is not entirely a
fabrication. The first is perfectly plausible; the second is what most of the
\emph{alien bacteria} commentary has done, and it is the version I would bet on. It is the
third possibility that I refuse to close the file on, because nothing else in the global
record asks me to.

The hospital end of the case has the same awkward texture. Staff at the time denied
receiving any such creature, and the Army inquiry found the ``aliens'' seen at the hospital
to be an expectant couple with dwarfism. Then, on the thirtieth anniversary in January 2026,
a neurologist named \'Italo Venturelli went on the record in the Globo documentary
\emph{O Mist\'erio de Varginha} saying that a colleague had asked him to look at something
that had just come out of surgery: child-sized, white rather than green or brown, a
teardrop-shaped skull, lilac eyes. ``I looked at it, it looked at me, it looked out the
window and back at me.'' He says he stayed silent for thirty years because he expected to be
called insane, and spoke only after a serious illness. Forensic pathologist Badan Palhares
--- the man who examined Josef Mengele's remains --- has confirmed only that the Army told him
material from Varginha would be sent to Campinas for analysis, which is a statement about a
phone call, not about a corpse. I have no way to test any of this, and a testimony first
given three decades after the event, into a camera, during an anniversary, is the weakest
class of evidence this book recognises. I include it because leaving it out would be
dishonest, not because it moves the needle.

Where that leaves me: the identification of the creature is unresolved, the cover-up claim
is unsupported, and the town has every commercial reason to keep the story warm --- there is
a water tower dressed as a flying saucer and an ET Memorial. And a young police officer is
still dead of something his doctors could not comfortably explain, on the timetable the
story requires. Varginha is the best example in the book of a case that survives not because
the evidence is good but because the ordinary explanation has one loose end it cannot tuck
in.

\subsubsection*{The rest of the world}

The 1977 events around Colares in Par\'a state involved dozens of villagers
reporting burns and puncture wounds from beams of light, sufficiently many that the
Brazilian Air Force ran a formal investigation --- Opera\c{c}\~ao Prato --- whose files
were later partly released. Whatever was happening at Colares, it left marks on people,
which makes it a different category of evidence from a light in the sky.

Belgium produced the best-documented mass event in European history. Between November
1989 and April 1990, thousands of witnesses across Wallonia reported large silent
triangular craft, and the Belgian Air Force cooperated openly with civilian
investigators --- an almost unique instance of a military organisation treating the
subject as a legitimate question rather than a security embarrassment. The famous radar
locks from the F-16 intercept of 30--31 March 1990 have since been substantially
reinterpreted, and the single best-known photograph from the wave was admitted as a hoax
by its author in 2011. That admission is why the Belgian wave is worth studying rather
than why it should be dismissed: the field corrected itself in public.

France is the only nation that has maintained a continuous, publicly funded, publicly
accessible investigation. GEIPAN, which sits inside the national space agency CNES,
has catalogued reports since 1977 and publishes its case files online with
classifications from A (identified) through D (unexplained after full investigation).
Roughly three per cent of its cases end in category D. That number --- from a
government agency with no stake in the answer and every institutional reason to prefer a
tidy file --- is the single most quotable statistic in this entire field.

France also supplies the field's most-examined patch of soil. On 8~January 1981 at
Trans-en-Provence, a retired technician named Renato Nicolai reported a dull metallic disc
settling briefly on his terrace and leaving behind a ring of compacted, scorched ground.
Gendarmes took soil and plant samples within a day, GEPAN --- GEIPAN's predecessor ---
analysed them in laboratories, and the results were genuinely odd: compression and heating
consistent with a heavy warm object, and chlorophyll depletion in nearby wild alfalfa that
fell off with distance from the trace. It remains the only case where a state agency has
published a physical dataset it could not account for. It also remains a single witness with
a single patch of ground, and the biochemical work has been contested ever since.

Britain's contribution is Rendlesham Forest in Suffolk, 26--28~December 1980, between the
twin USAF-operated bases at Bentwaters and Woodbridge. Airmen investigating lights in the
trees reported a small structured object; the deputy base commander, Lieutenant Colonel
Charles Halt, went out on the second night and wrote it up in a one-page memorandum to the
Ministry of Defence in January 1981, including elevated radiation readings at the landing
site. The memo is real, declassified, and does not say what people quote it as saying: it is
a careful, hedged, slightly baffled report by a serving officer. The candidate mundane
explanations --- the Orfordness lighthouse, a re-entering Cosmos rocket body, a bright star
low on the horizon --- account for a good deal of it and not obviously all of it.

Iran produced the single best-documented military intercept anywhere. On the night of
18--19~September 1976, two Imperial Iranian Air Force F-4 Phantoms were scrambled over
Tehran towards a brilliant object. Both crews reported losing instrumentation and
communications as they closed, and regaining them on breaking off; Lieutenant --- later
General --- Parviz Jafari has told the story publicly for decades. The resulting Defense
Intelligence Agency report was distributed to the White House, the NSA, the CIA and the
Joint Chiefs, and its evaluation section calls the case an outstanding example involving
credible observers and multiple sensor types. It is exactly the sort of document people
assume does not exist.

Australia's best mass-witness case is Westall, on 6~April 1966, when something like two
hundred students and staff at Westall High School in Clayton South, Melbourne, watched a
grey-white object descend near a paddock, then leave. There are no photographs, the official
file is thin to the point of insult, and the witnesses have spent fifty years being told they
saw a HIBAL balloon array. Zimbabwe's is the Ariel School near Ruwa on 16~September 1994,
where roughly sixty children reported a craft and a figure by the fence during morning break.
The South African researcher Cynthia Hind arrived within days and did the one thing that
matters with child witnesses --- she separated them and had them draw what they saw before
they could compare notes. Seventy-one drawings survive, and the convergence is striking.
The children are adults now, they have been re-interviewed at length, including for a 2021
BBC piece, and most have not budged. Convergent testimony from children is powerful and it
is also the single easiest thing in the world to contaminate; both statements are true and
the case sits on the line between them.

\begin{funfact}
The cultural variation runs deeper than shapes. Japanese reports feature far fewer
occupants and far more object-only sightings. Scandinavian cases skew heavily towards
lights and ``ghost rockets'', a wave of which crossed Sweden in 1946 --- a full year
before Kenneth Arnold. Chinese reports cluster around airports and are frequently
resolved as drones or launch debris, partly because reporting channels are different.
The phenomenon, whatever it is, appears to speak with a local accent. That is either
telling you something about the phenomenon or something about the witnesses, and
deciding which is most of the work.
\end{funfact}

What should you take from the international picture? Mainly this: the raw fact of
sightings is global and continuous, going back long before 1947, and no serious person
disputes it. What varies by culture is the interpretation --- the shapes, the occupants,
the motives, the cover-up. Section~\ref{sec:fallacylanguage} told you that language shapes the question. The
global record demonstrates that it also shapes the answer, which is why the only figures
I trust in this field are the ones produced by people who published their category~A
cases with the same enthusiasm as their category~D.

\begin{srcnotes}
\item Varginha chronology, witness ages and the \emph{Mudinho} finding per the Brazilian
Army's Military Police Inquiry (released via the Superior Military Court, digitised 2026)
and contemporaneous reporting, including Moffett, M., \emph{The Wall Street Journal},
28~June 1996.
\item Chereze's admission, surgery and death from generalised infection, and the statement
by former battalion commander Maur\'icio Ant\^onio Santos, per Brazilian press coverage and
the Army inquiry record.
\item Venturelli testimony per \emph{O Mist\'erio de Varginha} (Globo, 2026), reported in
\emph{El Pa\'is}, 20~January 2026. Sceptical summary per Dunning, B., \emph{Skeptoid}.
\item Trans-en-Provence: GEPAN/CNES Technical Note~16 (1983) and subsequent published
analyses of the soil and vegetation samples.
\item Rendlesham: Halt, C.\ ``Unexplained Lights'', memorandum to RAF Bentwaters/MoD,
13~January 1981, declassified.
\item Tehran: Defense Intelligence Agency report DIA 6~846~0139~76, 22~September 1976.
\item Ariel School: Hind, C., contemporaneous field notes and pupil drawings, September
1994; BBC coverage, 2021.
\end{srcnotes}
